% !TEX root = chapter-standalone.tex

\section{Coalgebras of a comonad}
\label{sec:coalgebras}


In \cref{sec:Eilenberg-Moore}, we saw that an algebra for a monad captures the idea of ``evaluating'' or ``collapsing'' formal expressions into concrete values: an algebra structure map $\monact \colon \monA(\Obja) \mto \Obja$ specifies how to reduce wrapped-up data. Dualizing this notion leads to co-algebras for a co-monad, which capture the idea of ``observing'' or ``decomposing'' a value into structured context: a co-algebra structure map $\comoncoact \colon \Obja \mto \comonW(\Obja)$ specifies how to produce an enriched view of an object.

\begin{ctdefinition}[Coalgebra of a comonad]
    \SYNDEF{coalgebra of a comonad}
    \label{def:comonad-coalgebra}
    Let~$\tup{\comonW, \comonex, \comondu}$ be a co-monad on a category~\CatC.
    A co-algebra of~$\comonW$ (also called a $\comonW$-co-algebra) is specified by: \

    \constit
    \begin{enumerate}
        \item an object $\Obja$ of~\CatC;
        \item a morphism $\comoncoact \colon \Obja \mto \comonW(\Obja)$ of \CatC.
    \end{enumerate}
    \condit
    \begin{enumerate}
        \item \emph{Cocomposition}: the following diagram commutes:
\begin{equation}\label{eq:coalg-cocomposition}
\begin{tikzcd}[every arrow/.append style={arrowmorphismstyle}]
    \Obja \arrow[r, "\comoncoact"] \arrow[d, "\comoncoact"'] & \comonW(\Obja) \arrow[d, "\comonW(\comoncoact)"] \\
    \comonW(\Obja) \arrow[r, "\comondu_\Obja"'] & (\comonWW)(\Obja)
\end{tikzcd}
\end{equation}
        \item \emph{Counit}: the following diagram commutes:
\begin{equation}\label{eq:coalg-counit}
\begin{tikzcd}[every arrow/.append style={arrowmorphismstyle}]
    \Obja \arrow[r, "\comoncoact"] \arrow[dr, "\catid_\Obja"'] & \comonW(\Obja) \arrow[d, "\comonex_\Obja"] \\
    & \Obja
\end{tikzcd}
\end{equation}

    \end{enumerate}
\end{ctdefinition}

\begin{remark}
\label{rem:coalgebra-duality}
One may observe that this definition is dual to the definition of an algebra for a monad (\cref{def:monad-algebra}): all the arrows in the diagrams are reversed.
More precisely, a $\comonW$-co-algebra in $\CatC$ is the same thing as a $\comonW$-algebra in $\CatC\op$ (where $\comonW$ is viewed as a monad on $\CatC\op$).
\end{remark}

\begin{remark}
\label{rem:coalgebra-conditions-spelled-out}
In terms of equations, the conditions say the following.
The cocomposition condition states that decomposing and then decomposing further (via $\comonW(\comoncoact)$) gives the same result as decomposing and then duplicating (via $\comondu_\Obja$):
\begin{equation}
\comoncoact \mthen \comonW(\comoncoact) = \comoncoact \mthen \comondu_\Obja.
\end{equation}
The counit condition states that extracting after decomposing gives back the original value:
\begin{equation}
\comoncoact \mthen \comonex_\Obja = \catid_\Obja.
\end{equation}
\end{remark}


\begin{example}[$\readerfunctor$-co-algebras for the product co-monad]
\label{exa:coalgebra-product-comonad}

Recall from \cref{exa:product-comonad} that the product co-monad is given by $\readerfunctor(\setA) = \setA \cartprod \subA$, with extraction $\comonex_\setA(\tup{\ela, \styleelements{s}}) = \ela$ and duplication $\comondu_\setA(\tup{\ela, \styleelements{s}}) = \tup{\tup{\ela, \styleelements{s}}, \styleelements{s}}$.

A $\readerfunctor$-co-algebra on a set $\setA$ is a function $\comoncoact \colon \setA \mto \setA \cartprod \subA$. We can write $\comoncoact(\ela) = \tup{\morb(\ela), \morc(\ela)}$ for some functions $\morb \colon \setA \mto \setA$ and $\morc \colon \setA \mto \subA$.

The counit condition requires $\comonex_\setA \mthen \comoncoact = \catid_\setA$, i.e.\ for all $\ela \setin \setA$:
\begin{equation*}
    \comonex_\setA(\comoncoact(\ela)) = \comonex_\setA(\tup{\morb(\ela), \morc(\ela)}) = \morb(\ela) = \ela.
\end{equation*}
So the counit condition forces $\morb = \catid_\setA$, meaning $\comoncoact(\ela) = \tup{\ela, \morc(\ela)}$.

The cocomposition condition requires $\comoncoact \mthen \comonW(\comoncoact) = \comoncoact \mthen \comondu_\setA$. Let us check this. On the one hand:
\begin{align*}
(\comoncoact \mthen \readerfunctor(\comoncoact))(\ela)
& = \readerfunctor(\comoncoact)(\tup{\ela, \morc(\ela)}) \\
& = \tup{\comoncoact(\ela), \morc(\ela)} \\
& = \tup{\tup{\ela, \morc(\ela)}, \morc(\ela)}.
\end{align*}
On the other hand:
\begin{align*}
(\comoncoact \mthen \comondu_\setA)(\ela)
& = \comondu_\setA(\tup{\ela, \morc(\ela)}) \\
& = \tup{\tup{\ela, \morc(\ela)}, \morc(\ela)}.
\end{align*}
These are always equal. So the cocomposition condition is automatically satisfied.

In summary, a $\readerfunctor$-co-algebra on $\setA$ is simply a function $\morc \colon \setA \mto \subA$, and the structure map is $\comoncoact(\ela) = \tup{\ela, \morc(\ela)}$.
\end{example}


\begin{example}[$\streamfun$-co-algebras for the stream co-monad]
\label{exa:coalgebra-stream-comonad}

Recall from \cref{exa:stream-comonad} that the stream co-monad is given by $\streamsof{\setA} = \setA^{\natnumbers}$, with extraction $\comonex_\setA(\styleelements{\sigma}) = \styleelements{\sigma}_0$ and duplication $\comondu_\setA(\styleelements{\sigma})_n = (\styleelements{\sigma}_n, \styleelements{\sigma}_{n+1}, \styleelements{\sigma}_{n+2}, \ldots)$.

A $\streamfun$-co-algebra on a set $\setA$ is a function $\comoncoact \colon \setA \mto \streamsof{\setA}$, that is, a function that assigns to each element $\ela \setin \setA$ an infinite stream $\comoncoact(\ela) \setin \setA^{\natnumbers}$.

The counit condition requires $\comoncoact \mthen \comonex_\setA = \catid_\setA$, i.e.\ for all $\ela \setin \setA$:
\begin{equation*}
    \comonex_\setA(\comoncoact(\ela)) = \comoncoact(\ela)_0 = \ela.
\end{equation*}
So the first element of the stream $\comoncoact(\ela)$ must be $\ela$ itself.

Now, define $\mora \colon \setA \mto \setA$ by $\mora(\ela) = \comoncoact(\ela)_1$; that is, $\mora$ extracts the second element of the stream. Given the counit condition, we have $\comoncoact(\ela)_0 = \ela$, and we claim that $\comoncoact(\ela)_n = \mora^n(\ela)$, where $\mora^n$ denotes $n$-fold composition of $\mora$.

To see this, we use the cocomposition condition, which requires $\comoncoact \mthen \streamfun(\comoncoact) = \comoncoact \mthen \comondu_\setA$.
Evaluating the right-hand side at position $n$:
\begin{align*}
(\comoncoact \mthen \comondu_\setA)(\ela)_n
& = \comondu_\setA(\comoncoact(\ela))_n \\
& = (\comoncoact(\ela)_n, \comoncoact(\ela)_{n+1}, \comoncoact(\ela)_{n+2}, \ldots).
\end{align*}
Evaluating the left-hand side at position $n$:
\begin{align*}
(\comoncoact \mthen \streamfun(\comoncoact))(\ela)_n
& = \comoncoact(\comoncoact(\ela)_n).
\end{align*}
So the cocomposition condition says that $\comoncoact(\comoncoact(\ela)_n) = (\comoncoact(\ela)_n, \comoncoact(\ela)_{n+1}, \ldots)$. In particular, looking at position $1$ of both sides: $\comoncoact(\comoncoact(\ela)_n)_1 = \comoncoact(\ela)_{n+1}$, which gives $\mora(\comoncoact(\ela)_n) = \comoncoact(\ela)_{n+1}$. By induction, starting from $\comoncoact(\ela)_0 = \ela$, we get $\comoncoact(\ela)_n = \mora^n(\ela)$.

In summary, a $\streamfun$-co-algebra on $\setA$ is the same thing as a function $\mora \colon \setA \mto \setA$, and the structure map sends each element to the stream of its iterates:
\begin{equation*}
    \comoncoact(\ela) = (\ela, \mora(\ela), \mora^2(\ela), \mora^3(\ela), \ldots).
\end{equation*}
In other words, $\streamfun$-co-algebras are \emph{deterministic dynamical systems}: an element of $\setA$ is a ``state'', and $\mora$ is the ``transition function'' that determines how the state evolves over time.
\end{example}


\begin{example}[$\stylefunctors{\aword{Store}}$-co-algebras for the store co-monad]
\label{exa:coalgebra-store-comonad}

Recall from \cref{exa:store-comonad} that the store co-monad is given by $\stylefunctors{\aword{Store}}(\setA) = \setA^{\subA} \cartprod \subA$, with extraction $\comonex_\setA(\morb, \styleelements{s}) = \morb(\styleelements{s})$ and duplication $\comondu_\setA(\morb, \styleelements{s}) = (\styleelements{s}' \mapsto (\morb, \styleelements{s}'),\; \styleelements{s})$.

A $\stylefunctors{\aword{Store}}$-co-algebra on a set $\setA$ is a function $\comoncoact \colon \setA \mto \setA^{\subA} \cartprod \subA$. We can write $\comoncoact(\ela) = (\morb_\ela, \styleelements{s}_\ela)$, where $\morb_\ela \colon \subA \mto \setA$ is a lookup table and $\styleelements{s}_\ela \setin \subA$ is a position.

The counit condition requires $\comoncoact \mthen \comonex_\setA = \catid_\setA$, i.e.\ for all $\ela \setin \setA$:
\begin{equation*}
    \comonex_\setA(\morb_\ela, \styleelements{s}_\ela) = \morb_\ela(\styleelements{s}_\ela) = \ela.
\end{equation*}
So the lookup table $\morb_\ela$, evaluated at the position $\styleelements{s}_\ela$, must return $\ela$ itself.

The cocomposition condition requires $\comoncoact \mthen \comondu_\setA = \comoncoact \mthen \stylefunctors{\aword{Store}}(\comoncoact)$. On the one hand:
\begin{align*}
(\comoncoact \mthen \comondu_\setA)(\ela)
& = \comondu_\setA(\morb_\ela, \styleelements{s}_\ela) \\
& = (\styleelements{s}' \mapsto (\morb_\ela, \styleelements{s}'),\; \styleelements{s}_\ela).
\end{align*}
On the other hand:
\begin{align*}
(\comoncoact \mthen \stylefunctors{\aword{Store}}(\comoncoact))(\ela)
& = \stylefunctors{\aword{Store}}(\comoncoact)(\morb_\ela, \styleelements{s}_\ela) \\
& = (\comoncoact \after \morb_\ela,\; \styleelements{s}_\ela) \\
& = (\styleelements{s}' \mapsto \comoncoact(\morb_\ela(\styleelements{s}')),\; \styleelements{s}_\ela).
\end{align*}
Comparing the first components at a position $\styleelements{s}' \setin \subA$: the left-hand side gives $(\morb_\ela, \styleelements{s}')$, and the right-hand side gives $\comoncoact(\morb_\ela(\styleelements{s}')) = (\morb_{\morb_\ela(\styleelements{s}')},\; \styleelements{s}_{\morb_\ela(\styleelements{s}')})$.

For these to be equal, we need two things. First, $\styleelements{s}_{\morb_\ela(\styleelements{s}')} = \styleelements{s}'$ for all $\styleelements{s}' \setin \subA$; in other words, the position assigned to any element in the range of $\morb_\ela$ must be the position where that element was looked up. Second, $\morb_{\morb_\ela(\styleelements{s}')} = \morb_\ela$ for all $\styleelements{s}'$; that is, every element looked up in the table carries the same table as the original.

In summary, a $\stylefunctors{\aword{Store}}$-co-algebra on $\setA$ is a pair of functions: a ``position'' function $\styleelements{s}_{(-)} \colon \setA \mto \subA$ and a ``context'' function $\morb_{(-)} \colon \setA \mto \setA^{\subA}$, such that looking up the position of $\ela$ in its own context returns $\ela$, and all elements in the same context share the same context and position assignment. This structure is closely related to the notion of a \emph{lens}: the position function ``gets'' a value from the state, and the context function provides a way to ``put'' any value back.
\end{example}


\begin{example}[$\stylefunctors{\aword{Cowrite}}$-co-algebras for the cowriter co-monad]
\label{exa:coalgebra-cowriter-comonad}

Recall from \cref{exa:cowriter-comonad} that the cowriter co-monad has $\stylefunctors{\aword{Cowrite}}(\setA) = \setA^{\monoidAset}$, with extraction $\comonex_\setA(\morb) = \morb(\idmon_\monoidA)$ and duplication $\comondu_\setA(\morb)(\styleelements{m}_1)(\styleelements{m}_2) = \morb(\styleelements{m}_1 \mtimesof\monoidA \styleelements{m}_2)$.

A $\stylefunctors{\aword{Cowrite}}$-co-algebra on a set $\setA$ is a function $\comoncoact \colon \setA \mto \setA^{\monoidAset}$. In other words, each element $\ela \setin \setA$ is assigned a function $\comoncoact(\ela) \colon \monoidAset \mto \setA$.

The counit condition requires $\comoncoact \mthen \comonex_\setA = \catid_\setA$, i.e.\ for all $\ela \setin \setA$:
\begin{equation*}
    \comonex_\setA(\comoncoact(\ela)) = \comoncoact(\ela)(\idmon_\monoidA) = \ela.
\end{equation*}
So the function $\comoncoact(\ela)$, evaluated at the identity element, must return $\ela$ itself.

The cocomposition condition requires $\comoncoact \mthen \comondu_\setA = \comoncoact \mthen \stylefunctors{\aword{Cowrite}}(\comoncoact)$. On the one hand:
\begin{align*}
(\comoncoact \mthen \comondu_\setA)(\ela)(\styleelements{m}_1)(\styleelements{m}_2)
& = \comondu_\setA(\comoncoact(\ela))(\styleelements{m}_1)(\styleelements{m}_2) \\
& = \comoncoact(\ela)(\styleelements{m}_1 \mtimesof\monoidA \styleelements{m}_2).
\end{align*}
On the other hand:
\begin{align*}
(\comoncoact \mthen \stylefunctors{\aword{Cowrite}}(\comoncoact))(\ela)(\styleelements{m}_1)(\styleelements{m}_2)
& = \comoncoact(\comoncoact(\ela)(\styleelements{m}_1))(\styleelements{m}_2).
\end{align*}
So the cocomposition condition says: $\comoncoact(\ela)(\styleelements{m}_1 \mtimesof\monoidA \styleelements{m}_2) = \comoncoact(\comoncoact(\ela)(\styleelements{m}_1))(\styleelements{m}_2)$.

If we write $\ela \cdot \styleelements{m} = \comoncoact(\ela)(\styleelements{m})$ for the ``action'' of $\monoidAset$ on $\setA$, the two conditions become:
\begin{align*}
    \ela \cdot \idmon_\monoidA & = \ela, \\
    \ela \cdot (\styleelements{m}_1 \mtimesof\monoidA \styleelements{m}_2) & = (\ela \cdot \styleelements{m}_1) \cdot \styleelements{m}_2.
\end{align*}
In other words, a $\stylefunctors{\aword{Cowrite}}$-co-algebra is precisely a \emph{right action} of the monoid $\monoidA$ on the set $\setA$.
\end{example}


\begin{ctdefinition}[$\comonW$-co-algebra morphism]
    \label{def:coalgebramorphism}
    Let~$\tup{\comonW, \comonex, \comondu}$ be a co-monad on a category~\CatC, and let~$\tup{\Obja_1, \comoncoact_1}$ and~$\tup{\Obja_2, \comoncoact_2}$ be co-algebras of $\comonW$.
    A morphism $\tup{\Obja_1, \comoncoact_1} \mto \tup{\Obja_2, \comoncoact_2}$ of $\comonW$-co-algebras is: \

    \constit
    \begin{enumerate}
        \item A morphism $\mora \colon \Obja_1 \mto \Obja_2$ in \CatC.
    \end{enumerate}
    \condit
    \begin{enumerate}
        \item The following diagram commutes:
\begin{equation}\label{eq:coalg-morph-compat}
\begin{tikzcd}[every arrow/.append style={arrowmorphismstyle}]
    \Obja_1 \arrow[r, "\mora"] \arrow[d, "\comoncoact_1"'] & \Obja_2 \arrow[d, "\comoncoact_2"] \\
    \comonW(\Obja_1) \arrow[r, "\comonW(\mora)"'] & \comonW(\Obja_2)
\end{tikzcd}
\end{equation}

    \end{enumerate}
\end{ctdefinition}

\begin{ctdefinition}
    [Category of $\comonW$-co-algebras]
    \label{def:catofcomonadcoalgebras}
    Let $\tup{\comonW, \comonex, \comondu}$ be a co-monad on a category \CatC.
    The \emph{category of $\comonW$-co-algebras} $\CatC_\comonW$ of the co-monad $\comonW$ is specified by:
    \begin{enumerate}
        \item \emph{Objects}: $\comonW$-co-algebras;
        \item \emph{Morphisms:} $\comonW$-co-algebra morphisms;
        \item \emph{Identities}: given a $\comonW$-co-algebra $\tup{\Obja, \comoncoact}$, its \SY{identity morphism} is $\catidat\Obja$;
        \item \emph{Composition}: is induced by the composition of morphisms in \CatC.
    \end{enumerate}
\end{ctdefinition}


\begin{gradedexercise}[\exname{CofreeCoalgebras}]
    \label{ex:CofreeCoalgebras}
    Let $\tup{\comonW, \comonex, \comondu}$ be a co-monad on a category $\CatC$, and let $\Obja$ be an object of $\CatC$.
    Your task is to prove that the object $\comonW(\Obja)$, together with the morphism
    \begin{equation}
        \comondu_\Obja \colon \comonW(\Obja) \mto (\comonWW)(\Obja)
    \end{equation}
    defines a co-algebra for the co-monad $\comonW$.
    \emph{Hint}: use the axioms/conditions in the definition of a co-monad.
\end{gradedexercise}
\solutionof{CofreeCoalgebras}

\begin{gradedexercise}[\exname{ProductComonadCoalgebraMorphisms}]
    \label{ex:ProductComonadCoalgebraMorphisms}
    Consider the product co-monad from \cref{exa:product-comonad} with $\readerfunctor(\setA) = \setA \cartprod \subA$.
    In \cref{exa:coalgebra-product-comonad} we saw that a $\readerfunctor$-co-algebra on a set $\setA$ is given by a function $\morc_\setA \colon \setA \mto \subA$.
    \begin{enumerate}
        \item Given $\readerfunctor$-co-algebras $\tup{\setA, \morc_\setA}$ and $\tup{\setB, \morc_\setB}$, spell out explicitly what the condition is for a function $\mora \colon \setA \mto \setB$ to be a morphism of $\readerfunctor$-co-algebras.
        \item Give a concrete example of sets $\setA$, $\setB$, $\subA$, functions $\morc_\setA$, $\morc_\setB$, and a co-algebra morphism between them.
    \end{enumerate}
\end{gradedexercise}
\solutionof{ProductComonadCoalgebraMorphisms}
